\scanpage{106}
\mspage{9}
d'où \deleted{en} en composant avec (8.2.6) un foncteur canonique
\[
 C^o\longrightarrow\underline{\mathrm{Hom}}\,(T(\widehat C),T(\mathrm{Ens}))\tag*{(8.2.7)}
\]
ou, de façon équivalente, un foncteur canonique
\[
 T(\widehat C)\xrightarrow{\alpha\,\simeq}\underline{\mathrm{Hom}}\,(C^o,S_T)\deleted{\unsure{(\mathrm{Ens})}}\quad\text{où }S_T=T(\mathrm{Ens})\tag*{(8.2.8)}
\]
Si on a un ensemble d'indices de base pour $T$, on aura un
carré commutatif
\diagram{\begin{tikzcd}[ampersand replacement=\&]
 {(8.2.9)}\& T(\widehat C)\arrow[r,"{\alpha\,\simeq}"]\arrow[d,"b_{\widehat C}"']\& \underline{\mathrm{Hom}}(C^o,S_T)\arrow[d,"{\underline{\mathrm{Hom}}(C^o,b_{(\mathrm{Ens})})}"]\\
 \&\widehat C{}^I\arrow[r,dash,"\simeq" description]\&\underline{\mathrm{Hom}}(C^o,(\mathrm{Ens}))^I
\end{tikzcd}}
\margin{$\begin{aligned}
\underline{\mathrm{Hom}}_\lambda(R,\widehat C)&\simeq\underline{\mathrm{Hom}}_{\lambda\text{ en }R}(R\times C^o,(\mathrm{Ens}))\\
&\simeq\underline{\mathrm{Hom}}_{\deleted{\unsure{\lambda}}}(\unsure{C^o|C},\underline{\mathrm{Hom}}_\lambda(R,(\mathrm{Ens})))
\end{aligned}$}

\deleted{\unsure{Lorsque}}\inserted{le foncteur} $\alpha$ (8.2.8) est \inserted{tjrs} une équivalence, \deleted{alors}\inserted{donc si $\underset{\rightarrow}{\lambda}=\varnothing$} le carré 2-cartésien
\deleted{\unsure{est-il}} équivalent à
\diagram{\begin{tikzcd}[ampersand replacement=\&]
 {(8.2.10)}\&T(C)\arrow[r,hook]\arrow[d,"b_C"]\&\underline{\mathrm{Hom}}(C^o,S_T)\arrow[d,"{\varphi\mapsto(b_{i\,\mathrm{Ens}}\circ\varphi)_{i\in I}}"]\\
 \&C^I\arrow[r,hook]\&\widehat C{}^I
\end{tikzcd}}
\margin{(2 cartésien \deleted{\illegible{3}} $(\underset{\rightarrow}{\lambda}=\varnothing)$ \deleted{\unsure{pour} \illegible{1}})}

\margin{$\underset{\rightarrow}{\lambda}=\varnothing$}
Ou, en termes imagés : la donnée, pour $(E_i)\in\mathrm{Ob}\,C^I$, d'une
structure d'espèce $T$ \inserted{dans $C$}, revient à la donnée \deleted{pour chaque}
kif-kif dans $\widehat C$ i.e. à la donnée pour tt $Z\in\mathrm{Ob}\,C$ d'une
structure d'espèce $T$ sur les \deleted{ensembles}\inserted{familles}
$(E_i(Z))_{i\inZ}=(\mathrm{Hom}(Z,E_i))_{i\in I}$
variant fonctoriellement avec $Z$.

-- Lorsque $\alpha$ est une \deleted{\unsure{isom}}\inserted{équivalence}, alors on voit donc que la cat
\inserted{(pour tt $C$ tel que $\widehat C$ une $\mathcal U$-catégorie)}
\deleted{structures} $T(C)$ est connue, à équivalence près (....), lorsque
l'on connaît
\[
\left\{\begin{aligned}
&S=S_T=T(\mathrm{Ens}),\quad\text{avec les foncteurs}\\
&\beta_i:S_T\longrightarrow(\mathrm{Ens})\qquad\beta=b_{i\,\mathrm{Ens}}
\end{aligned}\right.\tag*{(8.2.11)}
\]
comme la sous-catégorie pleine de $\underline{\mathrm{Hom}}(C^o,S_T)$ formée des
$\varphi:C^o\longrightarrow S_T$ tels que les $\beta_i\,S_T$ soient représentables.

8.2.12. Si on veut une car. intrinsèque, on peut dire aussi que
c'est la sous-catégorie pleine des $\varphi:C^o\longrightarrow S_T$ tels que
$\beta_{\mathrm{Ens}}\circ\varphi$ représentable pour tt $\beta\in R_T^\lambda$. Notons que le foncteur
\[
 \beta\longmapsto\beta_{\mathrm{Ens}}\ :\ R=R_T^\lambda\longrightarrow
 \underline{\mathrm{Hom}}(\underset{S_T}{\unsure{S|S_T}},(\mathrm{Ens}))=\widehat{S^o}\quad\text{cf. (3.5)}\tag*{(8.2.13)}
\]
est aussi le foncteur évident
\[
 R\longrightarrow\underline{\mathrm{Hom}}\,(\underline{\mathrm{Hom}}_\lambda(R,\mathrm{Ens}),(\mathrm{Ens}))\tag*{(8.2.14)}
\]
Or, si on considère $S=\underline{\mathrm{Hom}}_\lambda(R,(\mathrm{Ens}))$ comme ss-catégorie
pleine de $\widehat{R^o}$, \uline{celle-ci contient $R^o$} (\inserted{identifié à} la cat. des foncteurs
représentables de $R$ dans $(\mathrm{Ens})$) donc on a \deleted{$\unsure{R^o\subset S\subset\widehat{R^o}}$}
% archive sheet number: 105
